\documentclass[12pt]{article}

\usepackage[margin=1in]{geometry}
\usepackage{enumitem}
\usepackage{amsmath}
\usepackage{amssymb}
\usepackage{array}
\usepackage{titlesec}

\title{Lecture 12 Notes:\\ Aristotle's \textit{Topics}, Book V}
\author{}
\date{}

\begin{document}

\maketitle

\section*{Central Theme}

Aristotle's \textit{Topics}, Book V gives commonplace rules for testing
proposed properties.

\[
\boxed{
\textit{Topics}, \text{ Book V}
=
\text{commonplace rules for testing proposed properties}
}
\]

Book V teaches how to ask whether a proposed attribute really belongs uniquely
to a thing without being its definition.

\[
\boxed{
\text{A property is convertible with the subject, but does not state its essence.}
}
\]

A property must be clear, unambiguous, non-essential, non-repetitive, properly
qualified, and convertible with the subject. It must belong to the subject
alone, in the right respect, without merely naming the subject's essence.

\section*{1. From Genus to Property}

Book IV asked:

\[
\text{Is } G \text{ really the genus of } S?
\]

Book V asks:

\[
\text{Is } P \text{ really a property of } S?
\]

\subsection*{Whiteboard}

\[
\text{Genus} = \text{what the thing is broadly}
\]

\[
\text{Property} = \text{what belongs only to the thing, but does not state essence}
\]

\subsection*{Example}

\[
\text{man} \leftrightarrow \text{capable of receiving knowledge}
\]

But:

\[
\text{capable of receiving knowledge} \neq \text{definition of man}
\]

\subsection*{Teaching Point}

A property is convertible with the subject, but does not define it.

\[
\boxed{
\text{Property belongs uniquely without stating essence.}
}
\]

\section*{2. What a Property Is}

A property belongs to a subject alone and is predicated convertibly of it.

\[
S \Rightarrow P
\]

\[
P \Rightarrow S
\]

But a property does not state the essence.

\[
\text{Property} = \text{coextensive, non-essential predicate}
\]

\subsection*{Teaching Point}

A property follows the subject both ways.

\[
\boxed{
\text{Wherever the subject is, the property is; wherever the property is, the subject is.}
}
\]

\section*{3. Four Kinds of Property}

Aristotle distinguishes four kinds of property.

\[
\begin{array}{ll}
1. & \text{Essential property} \\
2. & \text{Relative property} \\
3. & \text{Permanent property} \\
4. & \text{Temporary property}
\end{array}
\]

\subsection*{Teaching Point}

Not every property is property in the same way.

\[
\boxed{
\text{Properties differ according to scope, relation, and duration.}
}
\]

\section*{4. Essential Property}

An essential property distinguishes its subject from everything else.

\subsection*{Example}

It is an essential property of man to be:

\[
\text{a mortal living being capable of receiving knowledge}
\]

This separates man from everything else.

\subsection*{Teaching Point}

An essential property is tested against all other things.

\[
\boxed{
\text{Essential property} = \text{distinguishing mark against everything}
}
\]

\section*{5. Relative Property}

A relative property distinguishes its subject not from everything, but from some
definite other thing.

\subsection*{Example}

The soul, in relation to the body, is fitted to command, while the body is
fitted to obey.

\[
\text{soul} : \text{body}
=
\text{command} : \text{obedience}
\]

\subsection*{Teaching Point}

A relative property is tested through comparison.

\[
\boxed{
\text{Relative property} = \text{distinguishing mark against a definite other}
}
\]

\section*{6. Permanent Property}

A permanent property belongs at every time and never fails.

\subsection*{Example}

It is a permanent property of God to be:

\[
\text{an immortal living being}
\]

\subsection*{Teaching Point}

A permanent property must hold across all times.

\[
\boxed{
\text{Permanent property} = \text{property always belonging}
}
\]

\section*{7. Temporary Property}

A temporary property belongs at a certain time.

\subsection*{Example}

It may be a temporary property of some particular man:

\[
\text{to be walking in the gymnasium}
\]

or:

\[
\text{to be walking now}
\]

\subsection*{Teaching Point}

A temporary property must be marked as temporary.

\[
\boxed{
\text{Temporary property} = \text{property at a specified time}
}
\]

\section*{8. Relative Properties Multiply the Questions}

A relative property can produce multiple lines of attack.

\subsection*{Two Problems}

Suppose someone says:

\[
\text{It is a property of man, in relation to horse, to be biped.}
\]

One may attack by showing:

\[
\text{man is not biped}
\]

or:

\[
\text{horse is biped}
\]

\subsection*{Four Problems}

Suppose someone says:

\[
\text{man is biped and horse is quadruped.}
\]

Then one may attack by showing:

\[
\text{man is not biped}
\]

\[
\text{man is quadruped}
\]

\[
\text{horse is biped}
\]

\[
\text{horse is not quadruped}
\]

\subsection*{Teaching Point}

Relative properties are arguable because they expose both sides of a
comparison.

\[
\boxed{
\text{A relative property gives the questioner more handles for attack.}
}
\]

\section*{9. The Most Arguable Properties}

Aristotle says the most arguable properties are:

\[
\begin{array}{ll}
1. & \text{essential properties} \\
2. & \text{permanent properties} \\
3. & \text{relative properties}
\end{array}
\]

An essential property can be tested against many things.

A permanent property can be tested across many times.

A relative property can be tested through several contrasted questions.

\subsection*{Teaching Point}

The more extensive the claim, the more dialectical pressure it can bear.

\[
\boxed{
\text{A wide property-claim creates many opportunities for argument.}
}
\]

\section*{10. Two Main Questions in Book V}

Book V asks two main questions.

\[
\boxed{
\text{Question 1: Has the property been correctly rendered?}
}
\]

\[
\boxed{
\text{Question 2: Is it really a property?}
}
\]

A statement may fail either because it is badly worded or because the proposed
attribute is not actually a property.

\subsection*{Teaching Point}

A property must be stated well before it can be tested well.

\[
\boxed{
\text{Bad formulation can ruin even a true insight.}
}
\]

\section*{11. A Property Should Be More Intelligible Than the Subject}

The terms in which the property is stated should be more intelligible than the
subject.

The purpose of giving a property is to make the subject better understood.

\subsection*{Bad Example}

It is not helpful to say that fire has as its property:

\[
\text{bearing a very close resemblance to the soul}
\]

because soul is less intelligible than fire.

\subsection*{Teaching Point}

A property should make the subject clearer, not more obscure.

\[
\boxed{
\text{Do not explain the clearer by the less clear.}
}
\]

\section*{12. The Attribution Must Also Be Intelligible}

It is not enough for the words themselves to be intelligible.

The attribution of the property to the subject must also be intelligible.

\subsection*{Bad Example}

Suppose someone says fire is:

\[
\text{the primary element wherein the soul is naturally found}
\]

The problem is not only the word ``soul.'' The whole attribution is obscure,
because it is unclear whether the soul is found in fire.

\subsection*{Teaching Point}

A property must be clear both in wording and in application.

\[
\boxed{
\text{Clear terms are not enough if the connection is unclear.}
}
\]

\section*{13. Avoid Ambiguity in the Property}

A property is badly rendered if the terms or the whole expression bear more than
one sense.

\subsection*{Example}

\[
\text{being sentient}
\]

may mean:

\[
\text{possessing sensation}
\]

or:

\[
\text{using sensation}
\]

If the distinction is not made, the property is obscure.

\subsection*{Teaching Point}

A property must be univocal enough to be argued about.

\[
\boxed{
\text{Ambiguity makes the property vulnerable to refutation in the wrong sense.}
}
\]

\section*{14. Avoid Ambiguity in the Subject}

A property can also fail if the subject itself is ambiguous.

\subsection*{Example}

\[
\text{the knowledge of this}
\]

may mean:

\[
\text{knowledge possessed by this}
\]

\[
\text{use of knowledge by this}
\]

\[
\text{knowledge about this}
\]

\[
\text{use of knowledge about this}
\]

No property can be clearly rendered unless the sense is distinguished.

\subsection*{Teaching Point}

You cannot assign a property clearly to an unclear subject.

\[
\boxed{
\text{Ambiguous subject, ambiguous property.}
}
\]

\section*{15. Avoid Repetition}

Aristotle warns against repeating the same term in the property.

This may happen directly or by replacing a term with an equivalent definition.

\subsection*{Direct Repetition}

\[
\text{the body which is the most rarefied of bodies}
\]

The word ``body'' has been repeated.

\subsection*{Repetition by Definition}

One may repeat the same term by replacing it with another expression that means
the same thing.

\subsection*{Teaching Point}

A property should be concise and non-circular in expression.

\[
\boxed{
\text{Repetition obscures rather than clarifies.}
}
\]

\section*{16. Avoid Universal Attributes}

A property should distinguish its subject.

Therefore, a property should not include terms common to everything.

\subsection*{Bad Example}

If someone defines a property of knowledge using:

\[
\text{unity}
\]

the rendering is suspect, because unity is a universal attribute.

\subsection*{Teaching Point}

A term that belongs to everything cannot distinguish the subject.

\[
\boxed{
\text{A property must distinguish.}
}
\]

\section*{17. Do Not Render Multiple Properties as One}

A property should not pile up several different properties unless this is
explicitly intended.

\subsection*{Bad Example}

\[
\text{fire is the most rarefied and lightest body}
\]

This gives more than one property.

\[
\text{most rarefied}
\]

\[
\text{lightest}
\]

Each may belong to fire alone, but the statement does not cleanly render one
property.

\subsection*{Teaching Point}

A property should state one distinguishing feature.

\[
\boxed{
\text{One property should not be disguised as a list.}
}
\]

\section*{18. Do Not Use the Subject or Its Species in the Property}

The property should make the subject more intelligible.

Therefore, it should not use the subject itself or one of its species.

\subsection*{Bad Example}

Do not state as a property of animal:

\[
\text{the substance to which man belongs as a species}
\]

because man is a species of animal and is posterior to animal.

\subsection*{Teaching Point}

A property should not explain the prior through the posterior.

\[
\boxed{
\text{Do not explain animal through man.}
}
\]

\section*{19. Avoid Opposites or Naturally Simultaneous Things}

A property should not be rendered through an opposite, something simultaneous by
nature, or something posterior to the subject.

\subsection*{Bad Example}

Do not render the property of good as:

\[
\text{the most direct opposite of evil}
\]

because evil is the opposite of good and is understood together with it.

\subsection*{Teaching Point}

A property should clarify from what is prior or better known.

\[
\boxed{
\text{Opposites and simultaneous correlates do not make the subject clearer.}
}
\]

\section*{20. A Permanent Property Should Always Follow}

A proposed permanent property fails if it sometimes ceases to be a property.

\subsection*{Bad Example}

It is not a good permanent property of animal:

\[
\text{sometimes to move and sometimes to stand still}
\]

because this does not always follow as a property.

\subsection*{Good Example}

It may be a property of virtue:

\[
\text{to make its possessor good}
\]

because this always follows virtue.

\subsection*{Teaching Point}

A permanent property must always accompany the subject.

\[
\boxed{
\text{What comes and goes is not a permanent property.}
}
\]

\section*{21. Temporary Properties Must Be Explicitly Marked}

If someone gives a property of the present time, he must say so.

\subsection*{Bad Example}

\[
\text{to be sitting with a particular man}
\]

This is temporary, and should not be presented as though it were permanent.

\subsection*{Good Example}

\[
\text{to be walking now}
\]

The temporal qualification is explicit.

\subsection*{Teaching Point}

Temporary properties require temporal qualification.

\[
\boxed{
\text{Say ``now'' when the property belongs only now.}
}
\]

\section*{22. Avoid Properties Known Only by Unstable Sensation}

A property should not depend on sensible evidence that can fail.

\subsection*{Bad Example}

Do not render the property of the sun as:

\[
\text{the brightest star that moves over the earth}
\]

because when the sun sets, sensation no longer verifies the movement.

\subsection*{Permissible Case}

A sensible attribute may be used when it necessarily belongs.

\[
\text{surface} = \text{the primary thing that is coloured}
\]

Here colour is sensible, but it belongs necessarily to surface in the relevant
respect.

\subsection*{Teaching Point}

A property should not depend on unstable perceptual access.

\[
\boxed{
\text{Sensation may help, but unstable sensation does not secure property.}
}
\]

\section*{23. Do Not Render the Definition as the Property}

A property should be convertible with the subject but should not state the
essence.

\[
\text{definition} = \text{convertible + essence}
\]

\[
\text{property} = \text{convertible + not essence}
\]

\subsection*{Bad Example}

\[
\text{man} = \text{walking biped animal}
\]

This is essence-like and therefore closer to definition than property.

\subsection*{Better Example}

\[
\text{man} = \text{naturally civilized animal}
\]

This may be convertible with man without strictly defining his essence.

\subsection*{Teaching Point}

Property must be convertible without being definition.

\[
\boxed{
\text{A property reveals a unique mark, not the essence itself.}
}
\]

\section*{24. Place the Subject Within Its Essence First}

Like definitions, properties should begin by placing the subject within its
essence, normally by giving the genus, then adding what distinguishes it.

\subsection*{Example}

A property of man may be rendered as:

\[
\text{animal capable of receiving knowledge}
\]

This first places man under:

\[
\text{animal}
\]

and then adds:

\[
\text{capable of receiving knowledge}
\]

\subsection*{Teaching Point}

A well-rendered property often imitates the order of definition without becoming
a definition.

\[
\boxed{
\text{Begin with the genus, then add the distinctive mark.}
}
\]

\section*{25. The Core Test: Convertibility}

After considering correct rendering, Aristotle turns to whether the proposed
attribute really is a property.

The core test is convertibility.

\[
S \Rightarrow P
\]

and:

\[
P \Rightarrow S
\]

The property must belong to every instance of the subject and only to that
subject.

\subsection*{Teaching Point}

A property follows the subject both ways.

\[
\boxed{
\text{Property} = \text{coextensive, non-essential predicate}
}
\]

\section*{26. The Property Must Belong to Every Subject in the Right Respect}

For destructive purposes, look whether the proposed property fails to belong to
any subject, or fails to belong in the relevant respect.

\subsection*{Example}

It is not a property of the man of science:

\[
\text{not to be deceived by an argument}
\]

because a geometer may be deceived if his figure is misdrawn.

\subsection*{Whiteboard}

\[
\text{every } S \text{ qua } S \Rightarrow P
\]

\subsection*{Teaching Point}

A property must belong to the subject in the respect in which it is named.

\[
\boxed{
\text{A property of } S \text{ must belong to } S \text{ qua } S.
}
\]

\section*{27. Test the Converse}

If the description applies where the name does not, the proposed property fails.

\subsection*{Bad Example}

\[
\text{living being that partakes of knowledge}
\]

is not a property of man, because it may apply to God, while the name ``man''
does not.

\[
P \text{ applies to non-}S
\Rightarrow
P \text{ is not property of } S
\]

\subsection*{Teaching Point}

A property must not be wider than its subject.

\[
\boxed{
\text{Where the description applies, the name must apply also.}
}
\]

\section*{28. Do Not Render the Subject as Property of Its Predicate}

Do not make the subject a property of something found in it.

\subsection*{Bad Example}

Do not render:

\[
\text{fire}
\]

as the property of:

\[
\text{the body with the most rarefied particles}
\]

This reverses subject and attribute.

\subsection*{Better Pattern}

Render what is found in the subject as the property of the subject.

\subsection*{Example}

\[
\text{specifically the heaviest body}
\]

may be rendered as a property of earth.

\subsection*{Teaching Point}

The property should belong to the subject, not reverse subject and attribute.

\[
\boxed{
\text{Do not make the subject into the property of its own predicate.}
}
\]

\section*{29. Do Not Render a Partaken Attribute as Property}

If the subject partakes of an attribute as part of its essence, that attribute
is not a property.

It is closer to a differentia.

\subsection*{Example}

If:

\[
\text{walking on two feet}
\]

is treated as partaken in the essence of man, it should not be rendered as a
property of man.

\subsection*{Teaching Point}

What belongs as essence is not property.

\[
\boxed{
\text{A property is not a part of the subject's essence.}
}
\]

\section*{30. Property Must Belong Simultaneously with the Subject}

A property should belong always and necessarily simultaneously with the subject,
unless it is explicitly temporary.

\[
S \Leftrightarrow P \text{ simultaneously}
\]

\subsection*{Bad Example}

\[
\text{walking through the market-place}
\]

may belong before or after the attribute ``man'' in the wrong way, and so it is
not a permanent property of man.

\subsection*{Good Example}

\[
\text{animal capable of receiving knowledge}
\]

belongs simultaneously and necessarily with man.

\subsection*{Teaching Point}

Property and subject must rise and fall together.

\[
\boxed{
\text{A property is co-present with its subject.}
}
\]

\section*{31. Same Things Should Have Same Properties}

If two terms mean the same thing, a property of one should be a property of the
other insofar as they are the same.

\[
S_1 = S_2
\Rightarrow
\text{same properties qua same}
\]

\subsection*{Example}

\[
\text{proper object of pursuit}
\]

and:

\[
\text{desirable}
\]

mean the same in the relevant respect.

If a proposed property fails for one, it fails for the other.

\subsection*{Teaching Point}

Synonymous subjects should carry the same properties in the same respect.

\[
\boxed{
\text{Sameness of subject requires sameness of property qua the same.}
}
\]

\section*{32. Same Kind Should Have Same-Kind Properties}

If subjects are the same in kind, their properties should be the same in kind.

\subsection*{Example}

Man and horse are the same in kind insofar as they are animals.

So proposed properties of each should be compared at the same kind-level.

\subsection*{Warning}

This rule can be deceptive if one property belongs to only one species while the
other belongs to many.

\subsection*{Teaching Point}

Properties should fit the kind-level at which the subject is considered.

\[
\boxed{
\text{Same-kind subjects call for same-kind properties.}
}
\]

\section*{33. Sophistical Difficulty: Subject Plus Accident}

A captious questioner may distinguish between a subject and the subject
qualified by an accident.

\subsection*{Example}

\[
\text{man}
\]

and:

\[
\text{white man}
\]

may be treated as different.

Then a property of man may appear to belong also to white man, and so seem not
to be unique.

\subsection*{Response}

The subject and subject-plus-accident are not absolutely different. They differ
in mode of being.

\subsection*{Teaching Point}

Do not let accidental qualification create fake counterexamples.

\[
\boxed{
\text{Man and white man differ accidentally, not absolutely.}
}
\]

\section*{34. Natural and Invariable Belonging}

Aristotle warns against confusing what belongs naturally with what belongs
invariably.

\[
\text{naturally belongs} \neq \text{invariably belongs}
\]

\subsection*{Example}

\[
\text{biped}
\]

may belong naturally to man, even though not every man has two feet at every
moment.

\subsection*{Teaching Point}

State whether the property belongs naturally, actually, or invariably.

\[
\boxed{
\text{Natural belonging must be marked as natural belonging.}
}
\]

\section*{35. Specify the Mode of Belonging}

Many errors arise because the speaker fails to say how the property belongs.

A property may belong:

\[
\begin{array}{ll}
1. & \text{naturally} \\
2. & \text{actually} \\
3. & \text{specifically} \\
4. & \text{absolutely} \\
5. & \text{primarily} \\
6. & \text{because the subject has a state} \\
7. & \text{because the subject is a state} \\
8. & \text{because the subject partakes of something} \\
9. & \text{because the subject is partaken of by something}
\end{array}
\]

\subsection*{Teaching Point}

The mode of belonging must be stated when ambiguity threatens.

\[
\boxed{
\text{A property claim must say not only what belongs, but how it belongs.}
}
\]

\section*{36. Natural, Actual, Specific, and Absolute Properties}

Some properties belong naturally, some actually, some specifically, and some
absolutely.

\subsection*{Natural}

\[
\text{biped belongs naturally to man}
\]

\subsection*{Actual}

\[
\text{having four fingers belongs actually to this particular man}
\]

\subsection*{Specific}

\[
\text{consisting of most rarefied particles belongs specifically to fire}
\]

\subsection*{Absolute}

\[
\text{life belongs absolutely to living being}
\]

\subsection*{Teaching Point}

Different modes of belonging require different formulations.

\[
\boxed{
\text{The same predicate may fail or succeed depending on how it is qualified.}
}
\]

\section*{37. State, Possessor of State, Partaking, and Being Partaken Of}

A property may also be stated in relation to a state or to partaking.

\subsection*{State and Possessor}

\[
\text{incontrovertible by argument}
\]

may be said of science as a state or of the scientist as the possessor of that
state, but these must not be confused.

\subsection*{Partaking}

\[
\text{sensation}
\]

may belong to animal, and to man because man partakes of animal.

\subsection*{Teaching Point}

State and possessor, partaker and partaken, must be distinguished.

\[
\boxed{
\text{Do not transfer a property across these relations without qualification.}
}
\]

\section*{38. Do Not State a Thing as Property of Itself}

A property cannot simply be the thing itself.

What shows the essence is definition, not property.

\subsection*{Bad Example}

If:

\[
\text{becoming}
\]

and:

\[
\text{beautiful}
\]

are taken as the same, then becoming cannot be a property of beautiful.

\subsection*{Better Example}

\[
\text{animate substance}
\]

may be a property of living creature if it is convertible without simply being
the same term.

\subsection*{Teaching Point}

A property must be convertible without being identical in essence.

\[
\boxed{
\text{Do not render the subject as its own property.}
}
\]

\section*{39. Things with Like Parts: Whole and Part Tests}

For things consisting of like parts, test whether the property of the whole is
true of the part, and whether the property of the part is true of the whole.

\subsection*{Bad Whole Example}

\[
\text{the sea} = \text{the largest volume of salt water}
\]

This may describe the whole sea, but not each part of the sea.

\subsection*{Bad Part Example}

\[
\text{air} = \text{breathable}
\]

Some air may be breathable, but the whole of air is not breathable.

\subsection*{Better Example}

Earth naturally falls downward, and this can be said of particular pieces of
earth taken collectively as earth.

\subsection*{Teaching Point}

For homogeneous things, test whether the property fits whole and part
appropriately.

\[
\boxed{
\text{A property of like-parted things must respect whole-part predication.}
}
\]

\section*{40. Use Contraries}

Aristotle then turns to opposites.

If the contrary of the proposed property is not a property of the contrary
subject, then the proposed property may fail.

\subsection*{Pattern}

\[
P \text{ property of } S
\]

\[
P' \text{ property of } S'
\]

where:

\[
P' = \text{contrary of } P
\]

\[
S' = \text{contrary of } S
\]

\subsection*{Example}

If desirable is a property of good, then objectionable may be a property of
evil.

\[
\text{good} \leftrightarrow \text{evil}
\]

\[
\text{desirable} \leftrightarrow \text{objectionable}
\]

\subsection*{Teaching Point}

Properties should correspond across contrary pairs.

\[
\boxed{
\text{Contrary subjects should carry contrary properties.}
}
\]

\section*{41. Use Relative Opposites}

Relative opposites also supply tests.

\subsection*{Example}

\[
\text{double} : \text{half}
\]

If it is a property of double to be in the proportion:

\[
2:1
\]

then it is a property of half to be in the proportion:

\[
1:2
\]

\subsection*{Teaching Point}

Properties of relatives should convert with their correlates.

\[
\boxed{
\text{Relative properties move with relative terms.}
}
\]

\section*{42. Use Privation and Possession}

If a property belongs to a state, the corresponding privative property should
belong to the privation, and vice versa.

\subsection*{Example}

If:

\[
\text{to see}
\]

is a property of:

\[
\text{sight}
\]

then:

\[
\text{failure to see}
\]

is a property of:

\[
\text{blindness}
\]

\subsection*{Teaching Point}

State and privation should carry corresponding properties.

\[
\boxed{
\text{Possession and privation mirror one another in property arguments.}
}
\]

\section*{43. Use Positive and Negative Terms}

Positive and negative terms provide another set of tests.

\subsection*{Example}

If:

\[
\text{to live}
\]

is a property of:

\[
\text{living being}
\]

then:

\[
\text{not to live}
\]

is a property of:

\[
\text{not-living being}
\]

\subsection*{Whiteboard}

\[
P \text{ property of } S
\Rightarrow
\neg P \text{ property of } \neg S
\]

\subsection*{Warning}

Some versions of this rule are deceptive.

A positive term does not belong to a negative subject at all, while a negative
term may belong to a positive subject without being its property.

\subsection*{Teaching Point}

Positive and negative subjects should be handled carefully.

\[
\boxed{
\text{Negation can guide property arguments, but it can also mislead.}
}
\]

\section*{44. Use Co-Ordinate Divisions}

If members of one division correspond to members of another, properties should
line up across those divisions.

\subsection*{Example}

If wisdom is:

\[
\text{the natural virtue of the rational faculty}
\]

then temperance may be:

\[
\text{the natural virtue of the faculty of desire}
\]

\subsection*{Teaching Point}

Parallel divisions allow parallel property arguments.

\[
\boxed{
\text{Co-ordinate divisions create co-ordinate properties.}
}
\]

\section*{45. Use Inflexions}

Aristotle says to look at inflected forms.

\subsection*{Example}

If:

\[
\text{beautifully}
\]

is not a property of:

\[
\text{justly}
\]

then:

\[
\text{beautiful}
\]

is not a property of:

\[
\text{just}
\]

\subsection*{Teaching Point}

A property claim should survive translation into related grammatical forms.

\[
\boxed{
\text{Test the adjective through the adverb, and the adverb through the adjective.}
}
\]

\section*{46. Use Inflexions of Opposites}

One should also examine the inflected forms of opposites.

\subsection*{Example}

If:

\[
\text{best}
\]

is a property of:

\[
\text{the good}
\]

then:

\[
\text{worst}
\]

may be a property of:

\[
\text{the evil}
\]

\subsection*{Teaching Point}

Opposite inflexions should correspond to opposite subjects.

\[
\boxed{
\text{The grammatical family and the opposite family should both be tested.}
}
\]

\section*{47. Use Likeness and Analogy}

Aristotle recommends looking at things in like relations.

\subsection*{Example}

\[
\text{builder}:\text{house}
=
\text{doctor}:\text{health}
\]

If producing health is not a property of doctor, then producing a house may not
be a property of builder.

\subsection*{Constructive Example}

\[
\text{doctor}:\text{ability to produce health}
=
\text{trainer}:\text{ability to produce vigor}
\]

If it is a property of a trainer to possess the ability to produce vigor, then
it may be a property of a doctor to possess the ability to produce health.

\subsection*{Teaching Point}

Analogy tests whether a property claim fits a parallel relation.

\[
\boxed{
\text{Parallel relations can support or undermine property claims.}
}
\]

\section*{48. Identically Related Predicates}

Aristotle also considers cases where the same predicate is identically related
to more than one subject.

This is mainly useful for destructive purposes, because the same thing cannot
usually be a property of more than one subject.

\subsection*{Example}

If prudence is knowledge of both the noble and the base, and if being knowledge
of the noble is not a property of prudence, then being knowledge of the base is
not a property either.

\subsection*{Teaching Point}

A property must be peculiar, not merely one of several identical relations.

\[
\boxed{
\text{The same predicate cannot usually be a property of several subjects.}
}
\]

\section*{49. Use Being, Becoming, and Destruction}

Property relations should remain coherent across being, becoming, and
destruction.

\subsection*{Pattern}

If:

\[
P \text{ is a property of } S
\]

in being, then corresponding forms may hold in:

\[
\text{becoming}
\]

and:

\[
\text{being destroyed}
\]

\subsection*{Example}

If being mortal is a property of man, then becoming mortal may be tied to
becoming man, and the destruction of mortal to the destruction of man.

\subsection*{Teaching Point}

Property relations should remain coherent across being, generation, and
destruction.

\[
\boxed{
\text{Being, becoming, and perishing should preserve the property relation.}
}
\]

\section*{50. Use the Idea of the Subject}

Aristotle includes a rule concerning the idea of the subject.

The proposed property must belong to the idea in the right respect.

\subsection*{Bad Example}

If being motionless belongs to man-itself only qua idea, and not qua man, then
motionless is not a property of man.

\subsection*{Good Example}

If living-creature-itself is compounded of soul and body qua living creature,
then this may support the property of living creature.

\subsection*{Teaching Point}

Even at the level of forms, the property must belong in the right respect.

\[
\boxed{
\text{The property must belong qua the subject, not merely qua something else.}
}
\]

\section*{51. Greater and Less}

Book V uses greater-and-less tests.

If more-P is not a property of more-S, then less-P is not a property of less-S,
nor P of S simply.

Conversely, if more-P is a property of more-S, then the lower and simple cases
may follow.

\subsection*{Pattern}

\[
\text{more } P \text{ property of more } S
\Rightarrow
P \text{ property of } S
\]

\subsection*{Example}

If higher sensation is a property of higher life, then sensation simply may be a
property of life simply.

\subsection*{Teaching Point}

Degree relations can confirm or destroy property relations.

\[
\boxed{
\text{More-and-less reasoning applies to properties as well as goods and genera.}
}
\]

\section*{52. Simple Predication and Degree Predication}

One may also move from simple property claims to qualified degrees.

\subsection*{Example}

If:

\[
\text{tendency to move upward by nature}
\]

is a property of fire, then:

\[
\text{greater tendency to move upward by nature}
\]

may be a property of what is more fiery.

\subsection*{Teaching Point}

If the simple property holds, the corresponding degree-property may hold in
degree cases.

\[
\boxed{
\text{Simple property can ground more-and-less property.}
}
\]

\section*{53. More Likely and Less Likely Properties}

If a more likely property fails to be a property of a more likely subject, then
a less likely property fails too.

\[
\text{more likely fails} \Rightarrow \text{less likely fails}
\]

If a less likely property succeeds, then a more likely property may succeed.

\[
\text{less likely succeeds} \Rightarrow \text{more likely succeeds}
\]

\subsection*{Example}

If perceiving is more likely to be a property of animal than knowing is of man,
and perceiving is not a property of animal, then knowing is not a property of
man.

\subsection*{Teaching Point}

Comparative plausibility supplies dialectical leverage.

\[
\boxed{
\text{The stronger case tests the weaker, and the weaker success supports the stronger.}
}
\]

\section*{54. Same Predicate, More Likely Subject}

If a predicate is more likely to be a property of one subject than another, and
it fails for the more likely subject, it also fails for the less likely.

\subsection*{Example}

To be coloured is more likely to be a property of surface than of body.

If it is not a property of surface, then it is not a property of body.

\subsection*{Teaching Point}

A predicate that fails where it most belongs cannot be secured where it less
belongs.

\[
\boxed{
\text{Failure in the stronger location refutes the weaker location.}
}
\]

\section*{55. More Likely Property of the Same Subject}

If a more likely property of a given subject fails, then the less likely
property fails as well.

\subsection*{Example}

Sensible is more likely than divisible to be a property of animal.

If sensible is not a property of animal, then divisible is not either.

\subsection*{Teaching Point}

Within the same subject, the more plausible property tests the less plausible.

\[
\boxed{
\text{If the better candidate fails, the worse candidate fails.}
}
\]

\section*{56. Like Degree Arguments}

If two candidate properties are equally likely, then establishing or destroying
one can affect the other.

\subsection*{Example}

It may be as much a property of the faculty of reason to be the primary seat of
wisdom as it is of the faculty of desire to be the primary seat of temperance.

If one is established, the other may be supported.

\subsection*{Teaching Point}

Like cases should be treated alike unless a difference is shown.

\[
\boxed{
\text{Equal plausibility creates parallel argumentative pressure.}
}
\]

\section*{57. Like Relations vs Like Belonging}

Aristotle distinguishes two kinds of comparison.

\subsection*{Like Relation}

This proceeds by analogy.

\[
A:B = C:D
\]

\subsection*{Like Belonging}

This proceeds by comparing whether attributes belong in an equal manner.

\[
P \text{ belongs to } S \text{ as much as } Q \text{ belongs to } T
\]

\subsection*{Teaching Point}

Analogy and equal belonging are related, but not identical.

\[
\boxed{
\text{Some comparisons are relational; others compare degrees of belonging.}
}
\]

\section*{58. Potential Properties and Nonexistent Correlates}

A potential property must be possible in relation to the correlate it requires.

\subsection*{Bad Example}

\[
\text{breathable}
\]

cannot be a property of air if it means:

\[
\text{able to be breathed}
\]

and no animals exist that can breathe it.

Air may exist without any breathing animal, but breathable air requires a
possible breather.

\subsection*{Teaching Point}

Potential properties must be possible in relation to their required correlates.

\[
\boxed{
\text{Check whether the potentiality has the correlate it needs.}
}
\]

\section*{59. Proper Potential Properties}

A potential property may succeed if it is rendered relative to something that
exists, or if the potentiality can belong even when its correlate does not
actually exist.

\subsection*{Example}

It may be a property of being:

\[
\text{capable of being acted upon or of acting}
\]

because being, when it exists, can be involved in acting or being acted upon.

\subsection*{Teaching Point}

Potential properties are valid when the relevant potentiality can genuinely
belong to the subject.

\[
\boxed{
\text{Potentiality must be grounded in the subject.}
}
\]

\section*{60. Avoid Superlatives}

Aristotle warns against stating a property in the superlative.

\subsection*{Bad Example}

\[
\text{fire} = \text{the lightest body}
\]

If fire perishes, there may still be some lightest body among the things that
remain.

So the description may survive where the name ``fire'' does not.

\subsection*{Teaching Point}

Superlative properties are unstable because the superlative may shift subjects.

\[
\boxed{
\text{A property should not depend on being the most or least among changing cases.}
}
\]

\section*{61. The Architecture of Book V}

Book V moves through the following arc:

\[
\text{what property is}
\rightarrow
\text{kinds of property}
\rightarrow
\text{correct rendering}
\rightarrow
\text{intelligibility and ambiguity}
\]

\[
\rightarrow
\text{avoid repetition and universals}
\rightarrow
\text{avoid essence and definition}
\rightarrow
\text{convertibility}
\rightarrow
\text{same subject and same kind}
\]

\[
\rightarrow
\text{natural, actual, temporary qualification}
\rightarrow
\text{whole and part}
\rightarrow
\text{opposites}
\rightarrow
\text{inflexions and analogy}
\]

\[
\rightarrow
\text{being, becoming, destruction}
\rightarrow
\text{greater and less}
\rightarrow
\text{potentiality and superlatives}
\]

\subsection*{Teaching Point}

Book V teaches how to test whether a proposed property is truly coextensive
with its subject without defining it.

\[
\boxed{
\text{The central question is whether the predicate belongs only to the subject, but not as essence.}
}
\]

\section*{62. Summary of Main Ideas}

\begin{enumerate}[itemsep=0.5em]
    \item Book V gives commonplace rules for testing proposed properties.

    \item A property belongs to the subject alone, convertibly, but without
    stating the essence.

    \item Properties may be essential, relative, permanent, or temporary.

    \item Relative properties generate multiple lines of attack.

    \item A property must be rendered in terms more intelligible than the
    subject.

    \item The attribution of the property to the subject must also be clear.

    \item Ambiguous terms and ambiguous subjects make properties badly rendered.

    \item A property should avoid repetition.

    \item A property should not use terms common to everything.

    \item A property should state one distinguishing feature, not a pile of
    several properties.

    \item A property should not use the subject itself or its species.

    \item A property should not be rendered through opposites or things
    simultaneous by nature.

    \item A permanent property must always follow the subject.

    \item A temporary property must be explicitly marked as temporary.

    \item A property should not depend on unstable sensation.

    \item A property must not be the definition.

    \item A well-rendered property often begins with the genus and then adds the
    distinctive mark.

    \item The central test of property is convertibility.

    \item A property must belong to every subject in the relevant respect.

    \item A property must not be wider than its subject.

    \item A property should not reverse subject and predicate.

    \item What belongs as essence is not property.

    \item Property and subject must belong simultaneously.

    \item Same subjects, insofar as they are the same, should have the same
    properties.

    \item Same-kind subjects should have same-kind properties.

    \item Accidental qualification should not create fake counterexamples.

    \item Natural, actual, specific, absolute, and temporary properties must be
    distinguished.

    \item A thing should not be rendered as a property of itself.

    \item Whole and part must be tested in things with like parts.

    \item Contraries, relatives, privation, and negation provide property tests.

    \item Co-ordinate divisions and inflexions help test properties.

    \item Analogy tests properties through parallel relations.

    \item Property relations should remain coherent across being, becoming, and
    destruction.

    \item Greater-and-less reasoning applies to properties.

    \item Potential properties must have the right relation to their correlates.

    \item Superlative properties are unstable and should be avoided.
\end{enumerate}

\section*{Final Framing}

\[
\boxed{
\textit{Topics}, \text{ Book V, teaches how to test whether a proposed property is truly coextensive with its subject without defining it.}
}
\]

Book V teaches that a property must be clear, unambiguous, non-essential,
non-repetitive, properly qualified, and convertible with the subject. It must
belong to the subject alone, in the right respect, without merely naming the
subject's essence.

\[
\boxed{
\text{A property is not what a thing is, but what belongs only to it.}
}
\]

\end{document}